% !TEX root = ../main.tex
\chapter{État de l'art et choix technologiques}\label{ch:etat-art}

\section*{Introduction}
Ce chapitre définit les concepts clés du projet, compare les solutions de prise de rendez-vous existantes, puis justifie les choix technologiques retenus.

\section{Concepts clés}
\begin{definition}{Logiciel en tant que service (SaaS)}
Modèle de distribution dans lequel l'application est hébergée par l'éditeur et utilisée par ses clients à travers le web, sur abonnement, sans installation locale.
\end{definition}

\begin{definition}{Multi-location}
Architecture dans laquelle une seule instance de l'application sert plusieurs clients (les «~locataires~», ici les cliniques), en isolant strictement leurs données.
\end{definition}

\subsection{Stratégies de multi-location}
Le tableau~\ref{tab:multi-location} compare les trois stratégies d'isolation des données les plus répandues.

\begin{table}[H]\centering
\caption{Stratégies d'isolation des données}\label{tab:multi-location}
\begin{tabularx}{\linewidth}{lXXl}
\toprule
\textbf{Stratégie} & \textbf{Avantages} & \textbf{Inconvénients} & \textbf{Coût} \\
\midrule
Base par locataire & Isolation maximale & Exploitation lourde & Élevé \\
Schéma par locataire & Bon compromis & Migrations multipliées & Moyen \\
Colonne \texttt{clinic\_id} & Simple, économique & Isolation à garantir dans le code & Faible \\
\bottomrule
\end{tabularx}
\end{table}

\section{Étude des solutions existantes}
Nous avons étudié trois solutions du marché selon des critères définis avec l'encadrant entreprise (tableau~\ref{tab:benchmark}). \todo{Remplacer « Solution A/B/C » par les vrais produits étudiés et citer leurs sites dans la bibliographie.}

\begin{table}[H]\centering
\caption{Comparaison des solutions existantes}\label{tab:benchmark}
\begin{tabularx}{\linewidth}{Xcccc}
\toprule
\textbf{Critère} & \textbf{Solution A} & \textbf{Solution B} & \textbf{Solution C} & \textbf{\nomprojet} \\
\midrule
Réservation en ligne 24 h/24 & \checkmark & \checkmark & \checkmark & \checkmark \\
Multi-cliniques & \checkmark & -- & \checkmark & \checkmark \\
Rappels SMS & \checkmark & \checkmark & -- & \checkmark \\
Interface en arabe & -- & -- & -- & \checkmark \\
Hébergement des données en Tunisie & -- & -- & -- & \checkmark \\
Tarif adapté au marché local & -- & \checkmark & -- & \checkmark \\
\bottomrule
\end{tabularx}
\end{table}

\section{Choix technologiques}
Le tableau~\ref{tab:choix} résume les technologies retenues et la raison de chaque choix.

\begin{table}[H]\centering
\caption{Choix technologiques}\label{tab:choix}
\begin{tabularx}{\linewidth}{llX}
\toprule
\textbf{Couche} & \textbf{Technologie} & \textbf{Justification} \\
\midrule
Back-end & Spring Boot 3 (Java 21) & Maîtrise de l'équipe, écosystème mature, sécurité intégrée~\cite{springboot}. \\
Front-end & React 18 + TypeScript & Composants réutilisables, typage statique~\cite{react}. \\
Base de données & PostgreSQL 16 & Robustesse, contraintes d'intégrité, index partiels~\cite{postgresql}. \\
Authentification & Keycloak & OIDC standard, gestion des rôles, fédération d'identités. \\
Messagerie & RabbitMQ & Envoi asynchrone des rappels, reprise en cas d'échec. \\
Conteneurs & Docker, Kubernetes & Déploiement reproductible, montée en charge~\cite{kubernetes}. \\
CI/CD & GitHub Actions & Intégré au dépôt, gratuit pour le projet. \\
\bottomrule
\end{tabularx}
\end{table}

\section*{Conclusion}
Ce chapitre a défini les concepts du projet, comparé les solutions du marché et justifié nos choix technologiques. Aucune des solutions étudiées ne réunit le multi-cliniques, l'arabe et l'hébergement local. Le chapitre suivant spécifie les besoins de \nomprojet.
